\documentclass[12pt]{article}
\usepackage[T1]{fontenc}
\usepackage[utf8]{inputenc}
\usepackage{lmodern}
\usepackage{textcomp}
\usepackage{enumitem}
\usepackage{geometry}
\geometry{margin=1in}
\begin{document}
\begin{center}
Answer Key - Biology - Undergraduate Level Assessment 23 \
\small (Answers are ordered exactly as in the question paper.)
\end{center}

\section*{Multiple Choice}
\begin{enumerate}[label=\arabic*.]
\item \textbf{Answer:} A
\\ \textit{Options:} A) They have flowers that bloom for several years before producing seeds; B) They have a unique leaf structure; C) They can survive in extremely cold climates; D) They produce large seeds; E) Their seeds can remain dormant for up to a century before germinating
\\ \textit{Note:} The correct answer highlights the unique reproductive strategy of cycads, as their prolonged flowering period before seed production distinguishes them from other seed plants, which typically have a more immediate reproductive cycle.
\item \textbf{Answer:} A
\\ \textit{Options:} A) The lateral-line system in fishes helps with the digestion of food.; B) The lateral-line system in fishes is used for excreting waste materials.; C) The lateral-line system in fishes is involved in the secretion of hormones.; D) The lateral-line system in fishes is used for breathing.; E) The lateral-line system in fishes is used for vision.
\\ \textit{Note:} The lateral-line system in fishes is primarily involved in detecting water movements and vibrations, which aids in locating food and avoiding predators, thereby indirectly supporting the digestion of food by facilitating hunting and foraging behaviors.
\item \textbf{Answer:} B
\\ \textit{Options:} A) dog, robin, fish; B) robin, dog; C) lizard, frog; D) fish, robin; E) fish, frog
\\ \textit{Note:} The robin and dog are endothermic animals that can regulate their body temperature internally, allowing them to maintain a constant body temperature despite fluctuations in environmental temperatures, while the other options are ectothermic and rely on external heat sources.
\item \textbf{Answer:} A
\\ \textit{Options:} A) Earthworms use gills, frogs use both gills and lungs, fish use lungs, and humans use skin.; B) Earthworms and humans use lungs, frogs and fish use gills.; C) All use gills for obtaining oxygen.; D) Earthworms and humans use gills, frogs use skin, and fish use lungs.; E) Earthworms use lungs, frogs use gills, fish use skin, and humans use gills.
\\ \textit{Note:} incorrect because earthworms absorb oxygen through their skin, frogs have both gills for aquatic respiration and lungs for terrestrial breathing, fish primarily use gills for extracting oxygen from water, and humans rely on lungs for respiration.
\item \textbf{Answer:} A
\\ \textit{Options:} A) By assessing the fossil content found within the rock layers; B) By comparing the density of the rock to a known age database; C) By examining the color and texture of the rock; D) By counting the layers of the rock; E) By measuring the amount of carbon-14 remaining in the rock
\\ \textit{Note:} The correct answer highlights the importance of fossil content in sedimentary rock layers, as the presence of certain fossils can help correlate the age of the rock with known time periods in geological history, thus aiding in age estimation.
\end{enumerate}

\section*{True / False}
\begin{enumerate}[label=\arabic*.]
\setcounter{enumi}{5}
\item \textbf{Answer:} True
\\ \textit{Options:} A) True; B) False
\\ \textit{Note:} Insulin facilitates the conversion of glucose to glycogen by activating glycogen synthase, thereby promoting glycogen storage in the liver and muscle tissues.
\item \textbf{Answer:} True
\\ \textit{Options:} A) True; B) False
\\ \textit{Note:} Sexual dimorphism often arises from parallel selection because both sexes adapt to similar environmental pressures, leading to traits that enhance reproductive success and survival.
\item \textbf{Answer:} False
\\ \textit{Options:} A) True; B) False
\\ \textit{Note:} Cyclic GMP primarily acts as a second messenger involved in signaling pathways such as vasodilation, while calcium ions are primarily released through the action of other second messengers like inositol trisphosphate (IP 3) and diacylglycerol (DAG).
\item \textbf{Answer:} False
\\ \textit{Options:} A) True; B) False
\\ \textit{Note:} Complete equilibrium in a gene pool cannot be achieved because evolutionary forces such as mutation, genetic drift, gene flow, and natural selection will always introduce changes, regardless of the population's size or isolation.
\item \textbf{Answer:} True
\\ \textit{Options:} A) True; B) False
\\ \textit{Note:} The destruction of the ozone layer reduces the amount of ultraviolet radiation that is absorbed, leading to increased UV exposure on Earth, which can harm organisms, disrupt ecosystems, and negatively affect the growth rates of populations dependent on those ecosystems.
\end{enumerate}

\section*{Long Form Response}
\begin{enumerate}[label=\arabic*.]
\setcounter{enumi}{10}
\item \textbf{Answer:} Natural selection can lead to the predominance of drug-resistant HIV variants in a patient undergoing treatment with 3 TC due to the initial presence of a few resistant strains in the viral population. When treatment begins, the antiretroviral drug selectively inhibits the growth of non-resistant viruses, allowing those few resistant variants to proliferate. This process is exacerbated by factors such as high mutation rates of the HIV virus, which generate genetic diversity, and the incomplete suppression of viral replication, which can occur if the drug is not taken consistently. Consequently, the resistant variants become more frequent, ultimately leading to a predominance of drug-resistant strains within the viral population.
\\ \textit{Note:} correct because it accurately describes how natural selection favors the survival and proliferation of pre-existing drug-resistant HIV variants in the presence of 3 TC treatment, influenced by the virus's high mutation rates and the selective pressure exerted by the drug.
\item \textbf{Answer:} Selaginella exhibits several evolutionary advancements over ferns, particularly in its vascular structure, gametophyte development, and embryo formation. Unlike ferns, which possess only tracheids in their xylem, Selaginella features a more advanced vascular system with both vessels in the xylem and sieve elements in the phloem, enhancing its efficiency in water and nutrient transport. Additionally, Selaginella undergoes autospory, resulting in a reduced and independent gametophyte that allows it to thrive in various environments. Furthermore, the embryo of Selaginella develops without a suspensor, which is a significant adaptation that supports early growth without reliance on maternal tissue. These features contribute to Selaginella's ability to occupy diverse ecological niches and demonstrate its evolutionary progression within the plant kingdom.
\\ \textit{Note:} correct because it accurately identifies Selaginella's advanced vascular structure, which includes vessels and sieve elements, as well as its unique gametophyte development through autospory, illustrating how these features enhance its evolutionary fitness compared to ferns.
\end{enumerate}

\vfill
\noindent\textit{End of Answer Key}
\end{document}